\begin{enumerate}[label=\textbf{\ansnum.\arabic*}]
  \item The factor is $\tfrac{(1+0.6)\,1.5}{1.8}=1.333$: the ball gets $1.33$ times the speed of the foot. Foot $1.5$\,m/s: ball $2.0$\,m/s; foot $2.5$\,m/s: ball $3.33$\,m/s.
  \item $v_b=\sqrt{2\cdot0.5\cdot6}=2.45$\,m/s; the foot speed is $2.45/1.333=1.84$\,m/s (the program gives $1.84$ for $6$\,m).
  \item $v_b\propto\sqrt{a}$: the ball speed must grow by $\sqrt2=1.41$, and so must the foot speed (the same factor, because the relation between them does not depend on the turf).
  \item The distance is proportional to the square of the speed: $1.1^2=1.21$, i.e. $21\%$ too far. For a kick of $12$\,m that is $2.5$\,m. To keep the error of the distance within $\pm10\%$ the speed must be within about $\pm5\%$.
  \item Command $(3,0)$: $(3,0)$ gives $1.0$; $(2.5,0)$: $\exp(-0.25/0.25)=0.368$; $(3,0.5)$: $0.368$; $(2,1)$: error$^2=2$, $\exp(-8)=0.0003\approx0$. A speed error of $0.5$\,m/s and a sideways error of $0.5$\,m/s are punished equally.
  \item If the reward is chosen by ``is the foot touching the ball now?'', then in the moment after the ball has left the foot the answer is no, and the reward switches back to the approach term, which now rewards bringing the foot to a ball that is moving away; at the same time the ball-velocity reward is lost. The policy then gets a reward signal that jumps depending on a millisecond-long event; it learns to \emph{chase} or \emph{touch repeatedly}. A latched flag (``the first contact has happened in this episode'') gives a clean sequence: first approach, then the ball velocity.
  \item At $0.3$\,m use at most about $0.02$\,m (noise grows with distance, so it is not larger than the value at $0.5$\,m). The real noise of the pipeline is low at close range in the ideal case (a pixel is a few millimetres at $1$\,m), but the ball is then often \emph{not visible at all}: it is below the field of view, hidden by the feet, or the head is moving. So the policy must work with a \emph{remembered} ball position (the last detection and its age, moved with the odometry), and the training must contain dropouts of the ball observation. During the approach the camera looks at the ball; during the kick the head should stay still so that the last estimates are good.
  \item $e=\sqrt{h'/h}=\sqrt{0.36}=0.60$. A range of $\pm25\%$ around $0.6$, i.e. $0.45$ to $0.75$, covers the measurement error (the height of a rebound is hard to read from video) and the differences between balls and temperatures, while still being physically plausible. (This is a suggestion, not a rule.)
  \item For $a=0.5$: $1.06,\,1.68,\,2.12,\,2.60$\,m/s; for $a=0.8$: $1.34,\,2.12,\,2.68,\,3.29$\,m/s. The relative difference is the same ($\sqrt{0.8/0.5}=1.26$, i.e. $26\%$ more), but the \emph{absolute} difference grows with the distance: $0.28$\,m/s at $2$\,m and $0.69$\,m/s at $12$\,m. For the long kicks the surface matters most, and a controller with a fixed foot speed will miss by metres (the distance changes by the factor $1.6$).
  \item An example: (1) ball fixed, command forward at one speed: success = ball within $\pm10^{\circ}$ and robot upright $>90\%$; watch the fall rate. (2) ball position $\pm5$\,cm: success $>90\%$; watch the contact rate. (3) direction $\pm30^{\circ}$ then $\pm90^{\circ}$: success on direction error $<10^{\circ}$; watch the support-foot behaviour. (4) speeds from $1$ to $3.5$\,m/s: success on the distance error $<20\%$; watch the joint torques and the motor limits. (5) noise, delay and randomisation: success $>90\%$ of the former, with a drop of the success rate less than $10$ points; watch the gap between the actor and the critic values.
\end{enumerate}
